\documentclass[11pt]{amsart}
\usepackage{amsmath,amssymb,amsthm}
\usepackage[margin=1.15in]{geometry}
\usepackage{enumerate}
\usepackage[hidelinks]{hyperref}

\newtheorem{theorem}{Theorem}[section]
\newtheorem{proposition}[theorem]{Proposition}
\newtheorem{lemma}[theorem]{Lemma}
\theoremstyle{definition}
\newtheorem{example}[theorem]{Example}
\theoremstyle{remark}
\newtheorem{remark}[theorem]{Remark}

\newcommand{\doi}[1]{\href{https://doi.org/#1}{doi:\nolinkurl{#1}}}
\newcommand{\arxiv}[1]{\href{https://arxiv.org/abs/#1}{arXiv:\nolinkurl{#1}}}
\newcommand{\N}{\mathbb N}
\newcommand{\Z}{\mathbb Z}
\newcommand{\R}{\mathbb R}
\newcommand{\E}{\mathbb E}
\newcommand{\Prob}{\mathbb P}
\newcommand{\cE}{\mathcal E}
\newcommand{\cF}{\mathcal F}
\newcommand{\cB}{\mathcal B}
\newcommand{\cT}{\mathcal T}
\newcommand{\dtv}{d_{\mathrm{TV}}}
\newcommand{\one}{\mathbf 1}
\setlength{\emergencystretch}{2.5em}

\begin{document}

\title{Setwise convergence of shifted laws without total variation convergence}
\author{Yitzchak Shmalo}
\address{Einstein Institute of Mathematics, The Hebrew University of Jerusalem, Jerusalem, Israel}
\email{yspennstate@gmail.com}
\subjclass[2020]{60G10, 37A25, 60B10}
\keywords{setwise convergence, total variation, stationary process, exact coupling, tail $\sigma$-field, mixing}

\begin{abstract}
Thorisson asked whether setwise convergence of the shifted laws of a discrete-time process to the law of a stationary
process forces convergence in total variation, that is, agreement on the tail $\sigma$-field. We show that it does not.
If the law of the process has a density with respect to a stationary mixing law, the shifted laws converge setwise,
and their total variation distance to the limit converges to half the $L^1$ distance from $1$ of the conditional expectation of the density given the
tail $\sigma$-field. Orbits of an invertible mixing transformation started from a law with a nonconstant density with respect to the
invariant measure are therefore
counterexamples, and zero-entropy mixing processes give counterexamples with two states. Processes with independent
coordinates give none. For Thorisson's companion question on couplings, we show that a single coupling under which the
shifted processes asymptotically agree on every event is sufficient but not necessary: for the horocycle flow, no
coupling brings the processes together in probability.
\end{abstract}

\maketitle

\section{Introduction}\label{sec:intro}

Let $(E,\cE)$ be a measurable space, write $\N=\{0,1,2,\dots\}$, and let $X=(X_k)_{k\in\N}$ and $X'=(X'_k)_{k\in\N}$
be stochastic processes with values in $E$. For $n\in\N$ let $\theta_n X=(X_{n+k})_{k\in\N}$ be the process shifted by
$n$, and let $\theta_n$ also denote the shift map on the path space $E^\N$, whose coordinates we write
$x=(x_k)_{k\in\N}$. A set $A\in\cE^{\N}$ is a \emph{tail set} if
$A\in\theta_n^{-1}(\cE^{\N})$ for every $n$, and \emph{invariant} if $\theta_1^{-1}A=A$.

Thorisson's coupling theory \cite{Tho00} gives two classical equivalences, recalled in his list of open problems
\cite{Tho11}, \cite[Theorem~3.1]{ThoPre}. First, the shifted laws $\Prob(\theta_nX\in\cdot\,)$ and $\Prob(\theta_nX'\in\cdot\,)$
become close in total variation if and only if $X$ and $X'$ have the same distribution on the tail sets, if and only if
they admit a distributional exact coupling. Second, the Ces\`aro averages of the shifted laws become close in total
variation if and only if $X$ and $X'$ agree on the invariant sets, if and only if they admit a distributional
shift-coupling. In particular, if the Ces\`aro averages converge setwise to the law of $X'$, then $X'$ is stationary and
the convergence of the averages is in total variation \cite[Corollary~3.1]{ThoPre}. Thorisson asked whether the
analogous statement holds without averaging \cite{Tho11}, \cite[Problem~3.1]{ThoPre}. Suppose
\begin{equation}\label{eq:setwise}
\Prob(\theta_nX\in A)\to\Prob(X'\in A)\quad\text{as }n\to\infty,\ \text{for all }A\in\cE^{\N},
\end{equation}
that is, suppose the shifted laws converge \emph{setwise} to the law of $X'$. Then $X'$ is stationary. Does it follow that $\Prob(\theta_nX\in\cdot\,)\to\Prob(X'\in\cdot\,)$ in total variation?
Thorisson guessed that the answer is negative, and asked, in that case, for a coupling characterization of
\eqref{eq:setwise} \cite[Problem~3.2]{ThoPre}.

We confirm the guess. Throughout, for probability measures $\mu,\mu'$ on one measurable space,
$\dtv(\mu,\mu')=\sup_A|\mu(A)-\mu'(A)|$ is the total variation distance, and $\one_A$ or $\one\{\cdot\}$
denotes an indicator function.

\begin{theorem}\label{thm:main}
Let $(\Omega,\cB,\nu,T)$ be a measure-preserving system in which $T$ is a bijection with measurable inverse, and assume
that it is mixing: $\nu(T^{-n}C\cap D)\to\nu(C)\nu(D)$ for all $C,D\in\cB$. Let $f\ge0$ with $\int f\,d\nu=1$, let
$Y$ have law $f\,d\nu$ and $Y'$ have law $\nu$, and put $E=\Omega$ and
\[
X_k=T^kY,\qquad X'_k=T^kY'\qquad(k\in\N).
\]
Then $X'$ is stationary, \eqref{eq:setwise} holds, and for every $n\in\N$
\[
\dtv\bigl(\Prob(\theta_nX\in\cdot\,),\Prob(X'\in\cdot\,)\bigr)=\tfrac12\int_\Omega|f-1|\,d\nu .
\]
In particular, if $f$ is not $\nu$-almost everywhere equal to $1$, the convergence \eqref{eq:setwise} does not hold in
total variation.
\end{theorem}

For the Bernoulli shift, that is, the two-sided shift on $\{0,1\}^\Z$ with the fair coin measure, and $f(y)=2\cdot\one\{y_0=1\}$, where
$y=(y_j)_{j\in\Z}$, the distance is $1/2$ for every $n$, and Example~\ref{ex:bernoulli} gives a tail set that
separates $X$ from $X'$. The state space of
Theorem~\ref{thm:main} is the space $\Omega$ itself. A two-point state space suffices as well
(Proposition~\ref{prop:binary}): if $Q$ is the law of a stationary mixing $\{0,1\}$-valued process of zero entropy, then
for each value $a$ with $Q(x_0=a)>0$ the process conditioned on $x_0=a$ satisfies \eqref{eq:setwise} at constant
distance $1-Q(x_0=a)$. Such processes exist by a theorem of Creutz, Pavlov and Rodock \cite{CPR23}.

Theorem~\ref{thm:main} is a consequence of a criterion for laws that are absolutely continuous with respect to a
stationary law (Proposition~\ref{prop:criterion}). Let $Q$ be a stationary probability measure on $E^\N$, let $X'$ have law $Q$, let $h$ be a probability
density with respect to $Q$, let $X$ have law $h\,dQ$, and write $\cF_n=\sigma(x_k:k\ge n)$ and $\cT=\bigcap_n\cF_n$ for the tail $\sigma$-field. Then
\begin{equation}\label{eq:intro-criterion}
\dtv\bigl(\Prob(\theta_nX\in\cdot\,),Q\bigr)=\tfrac12\bigl\|\E_Q[h\mid\cF_n]-1\bigr\|_{L^1(Q)}
\ \longrightarrow\ \tfrac12\bigl\|\E_Q[h\mid\cT]-1\bigr\|_{L^1(Q)}\qquad(n\to\infty),
\end{equation}
and \eqref{eq:setwise} holds for every such density $h$ if and only if $Q$ is mixing. Thus, when $Q$ is mixing,
\eqref{eq:setwise} always holds, and the convergence is in total variation if and only if $\E_Q[h\mid\cT]=1$
$Q$-almost surely. For laws with a density this is Thorisson's tail equivalence in quantitative form.

A measure-preserving transformation $T$ of $(\Omega,\cB,\nu)$ is \emph{exact} if $\bigcap_nT^{-n}\cB$ is trivial. The
transfer operator $P$ of $T$ maps the density of a law to the density of its image under $T$; when $T$ is invertible,
$Pu=u\circ T^{-1}$. Lin \cite{Lin71} proved that $T$ is exact if and only if $\|P^nu\|_{L^1(\nu)}\to0$ for every
$u\in L^1(\nu)$ with $\int u\,d\nu=0$. Mixing is the weaker property that $P^nu\to0$ weakly in $L^1(\nu)$ for every
such $u$. An invertible $T$ has $\bigcap_nT^{-n}\cB=\cB$, so it is never exact on a nontrivial space, and
Theorem~\ref{thm:main} is the path-space form of this fact. Indeed, $T^nY$ has density $f\circ T^{-n}$, whose
$L^1(\nu)$-distance from $1$ equals $\int|f-1|\,d\nu$ for every $n$ because $T$ preserves $\nu$, while mixing gives
$\int(\one_C\circ T^n)f\,d\nu\to\nu(C)$ for every $C\in\cB$.

Zeraoulia \cite{Zer26} constructs a process on $[0,1]$ with independent coordinates and oscillatory densities of
the form $1+\varepsilon\sin(2\pi mx)$, and states that its shifted laws converge setwise to a stationary limit
while their total variation distance from that limit stays constant.
Proposition~\ref{prop:independent} rules out every such process: a process with independent coordinates that
satisfies \eqref{eq:setwise} converges in total variation. If the $k$th coordinate has density
$1+\varepsilon\sin(2\pi m_kx)$, where $\varepsilon\in(0,1]$ is fixed and the $m_k$ are positive integers tending to
infinity, the shifted laws converge to the uniform product law on every set depending on finitely many coordinates,
but \eqref{eq:setwise} fails on a set depending on infinitely many coordinates (Example~\ref{ex:oscillating}).

Section~\ref{sec:coupling} concerns the coupling question. Suppose that $E$ is Polish and that
\eqref{eq:setwise} holds. Then for each $n$
separately there is a coupling of $\theta_nX$ with $X'$, and along these couplings the probability that a fixed event
occurs for one process and not for the other tends to zero (Lemma~\ref{lem:family}). So a characterization by
couplings says something new only if it uses a single coupling of $X$ and $X'$. For a stationary process $X'$ with
values in $E$, call a coupling of $X$ and $X'$ \emph{eventwise} if
\begin{equation}\label{eq:eventwise}
\Prob(\theta_nX\in A,\ \theta_nX'\notin A)\to0\quad\text{for every }A\in\cE^{\N}.
\end{equation}
An eventwise coupling gives \eqref{eq:setwise}, and for a Polish state space it is the same as \eqref{eq:setwise}
together with convergence of $X_n$ and $X'_n$ to each other in probability (Lemma~\ref{lem:closeness}). Eventwise
couplings exist in the examples of Theorem~\ref{thm:main} built on two-sided shifts whose invariant measure is trivial on $\bigcap_m\sigma(y_j:j\ge m)$, for instance the
Bernoulli shift of Example~\ref{ex:bernoulli} (Proposition~\ref{prop:Ktail}, where the coupling comes from Goldstein's maximal coupling
theorem \cite{Gol79}). They need not exist in general.

\begin{theorem}\label{thm:horocycle}
Let $G=\mathrm{SL}(2,\R)$, let $\Gamma\subset G$ be a cocompact lattice, let $\Omega=\Gamma\backslash G$ with its
$G$-invariant probability measure $\nu$ and a metric $d$ compatible with its topology, and let $T(x)=xu_1$, where
$u_t=\left(\begin{smallmatrix}1&t\\0&1\end{smallmatrix}\right)$. Let $f$ be a probability density with respect to $\nu$
that is not $\nu$-almost everywhere equal to $1$, and define $X$ and $X'$ as in Theorem~\ref{thm:main}. Then
\eqref{eq:setwise} holds, but there is no coupling of $X$ and $X'$ under which $d(X_n,X'_n)\to0$ in probability. In
particular $X$ and $X'$ admit no eventwise coupling.
\end{theorem}

The proof uses mixing of the horocycle flow and Ratner's equidistribution theorem for unipotent flows \cite{Rat91}.
Theorem~\ref{thm:horocycle} also bears on the last question in Thorisson's list for this topic
\cite[Problem~3.3]{ThoPre}, which asks for a single-coupling characterization of convergence in distribution of
$\theta_nX$ to $X'$ and suggests the property that $\theta_nX$ and $\theta_nX'$ approach each other in a metric for the product
topology. A recent note \cite{WS26} gives a binary
process whose shifted laws converge in distribution to an i.i.d.\ law while no coupling has this property, even in
probability. In Theorem~\ref{thm:horocycle} the shifted laws converge setwise, and still no coupling has this property.
The general coupling characterizations asked for in Problems~3.2 and~3.3 of \cite{ThoPre} remain open.


\section{A criterion for absolutely continuous laws}\label{sec:criterion}

We work on the path space $(E^\N,\cE^\N)$ with coordinates $x=(x_k)_{k\in\N}$ and shift $\theta=\theta_1$, so that
$\theta_n=\theta^n$. A probability measure $Q$ on $E^\N$ is \emph{stationary} if $Q\circ\theta^{-1}=Q$, and a
stationary $Q$ is \emph{mixing} if $Q(\theta^{-n}A\cap B)\to Q(A)Q(B)$ for all $A,B\in\cE^\N$. This is the usual notion of
ergodic theory; stronger notions under the same name, with a supremum over events of the remote future, are also in
use. Let
\[
\cF_n=\sigma(x_k:k\ge n),\qquad \cT=\bigcap_{n\in\N}\cF_n .
\]
The tail sets defined in Section~\ref{sec:intro} are exactly the elements of $\cT$, by the following lemma.

\begin{lemma}\label{lem:preimage}
For every $n$, $\{\theta^{-n}A:A\in\cE^\N\}=\cF_n$.
\end{lemma}

\begin{proof}
The left side is a $\sigma$-field, since preimages commute with complements and countable unions. Each
$\theta^{-n}A$ lies in $\cF_n$ because $\theta^n$ is $\cF_n$-measurable. Conversely, the generators
$\{x_{n+j}\in B\}$ ($j\in\N$, $B\in\cE$) of $\cF_n$ equal $\theta^{-n}\{x_j\in B\}$.
\end{proof}

We also recall that if $Q$ is a probability measure and $g\in L^1(Q)$ satisfies $\int g\,dQ=0$, then
\begin{equation}\label{eq:halfnorm}
\sup_{C\in\mathcal G}\Bigl|\int_C g\,dQ\Bigr|=\tfrac12\|g\|_{L^1(Q)}
\end{equation}
for every $\sigma$-field $\mathcal G$ with respect to which $g$ is measurable; the supremum is attained at
$C=\{g>0\}$, because the positive and negative parts of $g$ have the same integral.

\begin{proposition}\label{prop:criterion}
Let $Q$ be a stationary probability measure on $E^\N$, let $h\ge0$ be $\cE^\N$-measurable with $\int h\,dQ=1$, and let
$X$ be a process with $\Prob(X\in\cdot\,)=h\,dQ$. Write $P_n=\Prob(\theta_nX\in\cdot\,)$.
\begin{enumerate}[(a)]
\item If $Q$ is mixing, then $P_n(A)\to Q(A)$ for every $A\in\cE^\N$.
\item Conversely, if $P_n(A)\to Q(A)$ for every $A\in\cE^\N$ and every such $h$, then $Q$ is mixing.
\item For every $n$,
\[
\dtv(P_n,Q)=\tfrac12\bigl\|\E_Q[h\mid\cF_n]-1\bigr\|_{L^1(Q)} .
\]
This quantity is nonincreasing in $n$ and converges to $\frac12\|\E_Q[h\mid\cT]-1\|_{L^1(Q)}$.
\end{enumerate}
In particular, if $Q$ is mixing, then $P_n\to Q$ setwise, and $P_n\to Q$ in total variation if and only if
$\E_Q[h\mid\cT]=1$ $Q$-almost surely.
\end{proposition}

\begin{proof}
For $A\in\cE^\N$ we have $P_n(A)=\int\one_A\circ\theta^n\,h\,dQ$, and by stationarity $Q(A)=\int\one_A\circ\theta^n\,dQ$.

(a) Mixing says that $\int(\one_A\circ\theta^n)\,\one_B\,dQ\to Q(A)\int\one_B\,dQ$ for every $B$. By linearity the same
holds with $\one_B$ replaced by any simple function $s$. Given $\varepsilon>0$ choose a simple $s$ with
$\|h-s\|_{L^1(Q)}<\varepsilon$. Since $|\one_A\circ\theta^n|\le1$,
\[
\Bigl|P_n(A)-Q(A)\Bigr|\le\Bigl|\int(\one_A\circ\theta^n)\,s\,dQ-Q(A)\int s\,dQ\Bigr|+2\varepsilon ,
\]
and the first term tends to $0$.

(b) If $Q(B)=0$ there is nothing to prove. Otherwise take $h=\one_B/Q(B)$; then $P_n(A)=Q(\theta^{-n}A\cap B)/Q(B)$, and
$P_n(A)\to Q(A)$ is the mixing relation for $A$ and $B$.

(c) By the first line of the proof and Lemma~\ref{lem:preimage},
\[
\dtv(P_n,Q)=\sup_{A}\Bigl|\int(\one_A\circ\theta^n)(h-1)\,dQ\Bigr|=\sup_{C\in\cF_n}\Bigl|\int_C(h-1)\,dQ\Bigr|
=\sup_{C\in\cF_n}\Bigl|\int_C\bigl(\E_Q[h\mid\cF_n]-1\bigr)\,dQ\Bigr| ,
\]
and \eqref{eq:halfnorm} with $\mathcal G=\cF_n$ and $g=\E_Q[h\mid\cF_n]-1$ gives the formula. Since
$\E_Q[h\mid\cF_{n+1}]-1=\E_Q\bigl[\E_Q[h\mid\cF_n]-1\bigm|\cF_{n+1}\bigr]$ and conditional expectation is a contraction
of $L^1(Q)$, the sequence is nonincreasing. The $\sigma$-fields $\cF_n$ decrease to $\cT$, so by the reverse martingale
convergence theorem $\E_Q[h\mid\cF_n]\to\E_Q[h\mid\cT]$ in $L^1(Q)$, which gives the limit. The last statement combines
(a) and (c).
\end{proof}

\begin{remark}\label{rem:tail}
Part (c) agrees with Thorisson's tail characterization: $\int_C(\E_Q[h\mid\cT]-1)\,dQ=\Prob(X\in C)-Q(C)$ for
$C\in\cT$, so $\E_Q[h\mid\cT]=1$ holds exactly when $X$ and a process with law $Q$ have the same distribution on the
tail sets.
\end{remark}

\section{Counterexamples}\label{sec:examples}

We first deduce Theorem~\ref{thm:main} from Proposition~\ref{prop:criterion}.

\begin{proof}[Proof of Theorem~\ref{thm:main}]
Let $\Psi:\Omega\to\Omega^\N$, $\Psi(y)=(T^ky)_{k\in\N}$. Each coordinate of $\Psi$ is measurable, so $\Psi$ is
measurable, and $x_0\circ\Psi=\mathrm{id}$, $\theta\circ\Psi=\Psi\circ T$. Let $Q=\Psi_*\nu$ be the law of $X'=\Psi(Y')$.
Then $Q(\theta^{-1}A)=\nu(T^{-1}\Psi^{-1}A)=\nu(\Psi^{-1}A)=Q(A)$, so $X'$ is stationary, and
\[
Q(\theta^{-n}A\cap B)=\nu\bigl(T^{-n}\Psi^{-1}A\cap\Psi^{-1}B\bigr)\to\nu(\Psi^{-1}A)\,\nu(\Psi^{-1}B)=Q(A)Q(B),
\]
so $Q$ is mixing. Put $h(x)=f(x_0)$. For $A\in\cE^\N$,
\[
\Prob(X\in A)=\int\one_A(\Psi(y))f(y)\,\nu(dy)=\int\one_A(\Psi(y))\,h(\Psi(y))\,\nu(dy)=\int_A h\,dQ,
\]
so $\Prob(X\in\cdot\,)=h\,dQ$, and \eqref{eq:setwise} follows from Proposition~\ref{prop:criterion}(a).

Fix $n$ and let $h_n(x)=f(T^{-n}x_n)$, which is $\cF_n$-measurable because $T^{-n}$ is measurable. For every
$y\in\Omega$ we have $h_n(\Psi(y))=f(T^{-n}T^ny)=f(y)=h(\Psi(y))$, hence $\int|h_n-h|\,dQ=0$ and
$\E_Q[h\mid\cF_n]=h$ $Q$-almost surely. Proposition~\ref{prop:criterion}(c) gives
$\dtv(\Prob(\theta_nX\in\cdot\,),Q)=\frac12\int|h-1|\,dQ=\frac12\int|f-1|\,d\nu$.
\end{proof}

The proof also exhibits the tail set on which $X$ and $X'$ differ. For $B\in\cB$ let
\[
A_B=\bigl\{x\in\Omega^\N: T^{-k}x_k\in B\text{ for all but finitely many }k\bigr\}
=\bigcup_{r\ge n}\ \bigcap_{k\ge r}\{T^{-k}x_k\in B\}\qquad(n\in\N),
\]
which lies in $\cF_n$ for every $n$, hence in $\cT$. On the path $\Psi(y)$ we have $T^{-k}x_k=y$ for all $k$, so
$\Prob(X\in A_B)=\int_Bf\,d\nu$ and $\Prob(X'\in A_B)=\nu(B)$. For $B=\{f>1\}$ these differ by
$\frac12\int|f-1|\,d\nu$.

\begin{example}[Bernoulli shift]\label{ex:bernoulli}
Let $\Omega=\{0,1\}^\Z$ with the product $\nu$ of fair coin measures, $(Ty)_j=y_{j+1}$, and $f=2\cdot\one\{y_0=1\}$.
Thus $Y$ is a two-sided sequence of fair bits conditioned on $Y_0=1$, and the state of $\theta_nX$ at time $k$ is the
sequence $(Y_{j+n+k})_{j\in\Z}$.

\emph{Setwise convergence.} Let $A\in\cE^\N$ and $C=\Psi^{-1}A$, so that $\Prob(\theta_nX\in A)=\Prob(T^nY\in C)$ and
$\Prob(X'\in A)=\nu(C)$. Given $\varepsilon>0$, choose a set $C'$ depending only on the coordinates in $[-m,m]$ for some
$m$, with $\nu(C\,\triangle\,C')<\varepsilon$, where $\triangle$ denotes symmetric difference; such sets form an
algebra generating $\cB$. The law of $T^nY$ is
$(f\circ T^{-n})\,d\nu$, with density at most $2$, so
$|\Prob(T^nY\in C)-\Prob(T^nY\in C')|\le2\varepsilon$. For $n>m$ the coordinates $(T^nY)_j=Y_{j+n}$, $|j|\le m$, do not
involve $Y_0$, so $\Prob(T^nY\in C')=\nu(C')$. Hence $|\Prob(\theta_nX\in A)-\Prob(X'\in A)|\le3\varepsilon$ for $n>m$.

\emph{No convergence in total variation.} The event $\{x:(x_0)_{-n}=1\}$ has probability $\Prob(Y_0=1)=1$ under
$\theta_nX$ and probability $1/2$ under $X'$, so the distance is at least $1/2$; it equals $1/2$ by
Theorem~\ref{thm:main}.

\emph{A separating tail set.} Let $A=\{x:(x_k)_{-k}=1\text{ for all but finitely many }k\}$, a tail set. Since
$(X_k)_{-k}=Y_0=1$ and $(X'_k)_{-k}=Y'_0$ for every $k$, we get $\Prob(X\in A)=1$ and $\Prob(X'\in A)=1/2$.
\end{example}

Next we show that the state space can be taken to have two points. Let $Q$ be a stationary probability measure on
$\{0,1\}^\N$ and $H_m=-\sum_{w\in\{0,1\}^m}Q(x_0\cdots x_{m-1}=w)\log Q(x_0\cdots x_{m-1}=w)$. The entropy of $Q$ is
$\lim_mH_m/m$.

\begin{lemma}\label{lem:zeroentropy}
If $Q$ is a stationary probability measure on $\{0,1\}^\N$ of zero entropy, then $x_0$ is $Q$-almost surely equal to an
$\cF_n$-measurable function, for every $n$.
\end{lemma}

\begin{proof}
Write $H(\,\cdot\mid\cdot\,)$ for conditional entropy under $Q$. By stationarity
$H(x_0\mid x_1,\dots,x_m)=H_{m+1}-H_m$, and this is nonincreasing in $m$ because conditioning on more variables does
not increase entropy. Since $H_m/m$ is the Ces\`aro mean of $H_1,\,H_2-H_1,\dots,H_m-H_{m-1}$, the limit of
$H_{m+1}-H_m$ is the entropy, which is $0$. Let $p_m(a)=Q(x_0=a\mid x_1,\dots,x_m)$ and
$p_\infty(a)=Q(x_0=a\mid\cF_1)$. By martingale convergence $p_m(a)\to p_\infty(a)$ almost surely. The function
$\ell(t)=-t\log t$, with $\ell(0)=0$, is bounded and continuous on $[0,1]$, so by bounded convergence
\[
0=\lim_m H(x_0\mid x_1,\dots,x_m)=\lim_m\E_Q\Bigl[\sum_a\ell(p_m(a))\Bigr]=\E_Q\Bigl[\sum_a\ell(p_\infty(a))\Bigr].
\]
Since $\ell\ge0$ on $[0,1]$ with equality only at $0$ and $1$, $p_\infty(a)\in\{0,1\}$ almost surely. Moreover
\[
Q\bigl(x_0=a,\ p_\infty(a)=0\bigr)=\E_Q\bigl[p_\infty(a)\,\one\{p_\infty(a)=0\}\bigr]=0 .
\]
Hence $x_0=\beta(x_1,x_2,\dots)$ almost surely, where the measurable function $\beta$ equals $a$ on the set
$\{p_\infty(a)=1\}$. By stationarity $x_j=\beta(x_{j+1},x_{j+2},\dots)$ almost surely for every $j$. Substituting
these identities for $j=1,\dots,n-1$ into one another shows that $x_0$ is almost surely a measurable function of
$(x_n,x_{n+1},\dots)$.
\end{proof}

\begin{proposition}\label{prop:binary}
Let $Q$ be a stationary mixing probability measure on $\{0,1\}^\N$ of zero entropy, let $a\in\{0,1\}$ with
$p=Q(x_0=a)\in(0,1)$, let $X'$ have law $Q$, and let $X$ have law $Q(\,\cdot\mid x_0=a)$. Then \eqref{eq:setwise}
holds, and $\dtv(\Prob(\theta_nX\in\cdot\,),Q)=1-p$ for every $n$. Such measures $Q$ exist.
\end{proposition}

\begin{proof}
The law of $X$ is $h\,dQ$ with $h=\one\{x_0=a\}/p$. By Lemma~\ref{lem:zeroentropy}, $h$ agrees $Q$-almost surely with
an $\cF_n$-measurable function, so $\E_Q[h\mid\cF_n]=h$, and Proposition~\ref{prop:criterion} gives
\eqref{eq:setwise} and
\begin{align*}
\dtv\bigl(\Prob(\theta_nX\in\cdot\,),Q\bigr)&=\tfrac12\|h-1\|_{L^1(Q)}
=\tfrac12\Bigl(Q(x_0=a)\bigl(\tfrac1p-1\bigr)+Q(x_0\ne a)\Bigr)\\
&=\tfrac12\bigl((1-p)+(1-p)\bigr)=1-p .
\end{align*}

For existence, Creutz, Pavlov and Rodock \cite{CPR23} construct rank-one transformations, the extremely elevated
staircases, that are mixing \cite[Theorem~3.1]{CPR23}. Each of them is
measure-theoretically isomorphic to the shift on a subshift of $\{0,1\}^\Z$ with an invariant probability measure
$\mu$. The word complexity of that subshift is the number $c(q)$ of words of length $q$ occurring in it, and it can be
made $o(q^2)$ as $q\to\infty$ \cite[Theorem~4.1]{CPR23}, applied with the function $q\mapsto q^2$. Let $Q$ be the law of the coordinates with nonnegative index under $\mu$. It is
stationary and mixing because $\mu$ is. Its entropy is $0$, since a block of length $q$ takes at most $c(q)$ values,
so that $H_q\le\log c(q)$, and $\log c(q)/q\to0$. Finally $x_0$ is not almost surely constant, since
otherwise every coordinate would be almost surely equal to the same constant and $\mu$ would be a point mass, whereas
$\mu$ is nonatomic.
\end{proof}

\begin{remark}\label{rem:position}
Convergence \eqref{eq:setwise} lies strictly between the two notions of Thorisson's theory. It implies setwise
convergence of the Ces\`aro averages, hence, by \cite[Corollary~3.1]{ThoPre}, agreement on invariant sets, distributional
shift-coupling and convergence of the Ces\`aro averages in total variation. Theorem~\ref{thm:main} shows that it does not
imply convergence in total variation. Conversely, convergence of the Ces\`aro averages in total variation does not imply
\eqref{eq:setwise}: let $X'$ be the stationary two-state chain that alternates deterministically between $0$ and $1$,
started from the uniform distribution, and let $X$ be the same chain started at $0$. Then $\Prob(\theta_nX\in\{x_0=0\})$
alternates between $1$ and $0$, while $\Prob(X'\in\{x_0=0\})=1/2$. More precisely, the law of $\theta_kX$ is the point
mass at the alternating path starting with $0$ for even $k$ and at the one starting with $1$ for odd $k$, and the law
of $X'$ gives mass $1/2$ to each of these two paths. So the average of the first $N$ shifted laws of
$X$ equals the law of $X'$ when $N$ is even, and differs from it by $1/(2N)$ in total variation when $N$ is odd.
\end{remark}

\section{Independent coordinates}\label{sec:independent}

For probability measures $\mu,\mu'$ on a measurable space, the Hellinger affinity is
$\rho(\mu,\mu')=\int\sqrt{pp'}\,d\lambda$, where $p,p'$ are densities with respect to any measure $\lambda$ dominating
both; it does not depend on $\lambda$. We use two standard facts. First, $\rho(\mu,\mu')=0$ if and only if $\mu$ and
$\mu'$ are mutually singular. Second, by the Cauchy--Schwarz inequality applied to
$|p-p'|=|\sqrt p-\sqrt{p'}|\,(\sqrt p+\sqrt{p'})$,
\begin{equation}\label{eq:lecam}
\dtv(\mu,\mu')=\tfrac12\int|p-p'|\,d\lambda\le\tfrac12\Bigl(\int(\sqrt p-\sqrt{p'})^2d\lambda\Bigr)^{1/2}
\Bigl(\int(\sqrt p+\sqrt{p'})^2d\lambda\Bigr)^{1/2}=\sqrt{1-\rho(\mu,\mu')^2}.
\end{equation}

Kakutani \cite{Kak48} proved the following product formula when $\mu_k$ and $\lambda_k$ are equivalent for every
$k$. It holds for arbitrary probability measures.

\begin{lemma}\label{lem:kakutani}
Let $\mu_k,\lambda_k$ $(k\in\N)$ be probability measures on $(E,\cE)$. Then
$\rho\bigl(\bigotimes_k\mu_k,\bigotimes_k\lambda_k\bigr)=\prod_k\rho(\mu_k,\lambda_k)$.
\end{lemma}

\begin{proof}
Let $M=\bigotimes_k\mu_k$, $L=\bigotimes_k\lambda_k$, $R=\frac12(M+L)$, and let $p=dM/dR$, $q=dL/dR$. Let $\mathcal
G_d=\sigma(x_0,\dots,x_{d-1})$ and $p_d=\E_R[p\mid\mathcal G_d]$, $q_d=\E_R[q\mid\mathcal G_d]$, which are the densities
of the restrictions of $M$ and $L$ to $\mathcal G_d$ with respect to that of $R$. The affinity of the restrictions is
$\int\sqrt{p_dq_d}\,dR=\prod_{k<d}\rho(\mu_k,\lambda_k)$, by Fubini's theorem with the dominating product measure
$\bigotimes_{k<d}\frac12(\mu_k+\lambda_k)$ and the independence of the affinity from the dominating measure. By
martingale convergence $p_d\to p$ and $q_d\to q$ in $L^1(R)$, and since $(\sqrt s-\sqrt t)^2\le|s-t|$ for $s,t\ge0$,
$\sqrt{p_d}\to\sqrt p$ and $\sqrt{q_d}\to\sqrt q$ in $L^2(R)$. Hence $\int\sqrt{p_dq_d}\,dR\to\int\sqrt{pq}\,dR$.
\end{proof}

\begin{proposition}\label{prop:independent}
Let $X$ have independent coordinates, $X_k$ with law $\mu_k$, and suppose \eqref{eq:setwise} holds for some process
$X'$. Then $X'$ is i.i.d., and $\Prob(\theta_nX\in\cdot\,)\to\Prob(X'\in\cdot\,)$ in total variation.
\end{proposition}

\begin{proof}
Taking $A=\{x_0\in B\}$ in \eqref{eq:setwise} gives $\mu_n(B)\to\lambda(B):=\Prob(X'_0\in B)$ for every $B\in\cE$, hence
$\mu_{n+j}(B)\to\lambda(B)$ for every fixed $j$. Taking $A=\{x_0\in B_0,\dots,x_d\in B_d\}$ gives
$\Prob(X'_0\in B_0,\dots,X'_d\in B_d)=\lim_n\prod_{j\le d}\mu_{n+j}(B_j)=\prod_{j\le d}\lambda(B_j)$. These sets form a
$\pi$-system generating $\cE^\N$, so $Q:=\Prob(X'\in\cdot\,)=\lambda^{\otimes\N}$. Let
$P_n=\Prob(\theta_nX\in\cdot\,)=\bigotimes_{j\ge0}\mu_{n+j}$, a measure on the same space $E^\N$ for every $n$,
and $r_k=\rho(\mu_k,\lambda)\in[0,1]$. By Lemma~\ref{lem:kakutani}, $\rho(P_n,Q)=\prod_{j\ge0}r_{n+j}=\prod_{k\ge
n}r_k$, which is nondecreasing in $n$ because the factors lie in $[0,1]$.

Suppose first that $\prod_{k\ge n}r_k=0$ for every $n$. Then $P_n$ and $Q$ are mutually singular for every $n$: there
are sets $S_n$ with $P_n(S_n)=1$ and $Q(S_n)=0$. The set $S=\bigcup_nS_n$ satisfies $Q(S)=0$ and $P_n(S)=1$ for every
$n$, contradicting \eqref{eq:setwise} for $A=S$. Hence $a:=\prod_{k\ge n_0}r_k>0$ for some $n_0$. Then $r_k>0$ for
$k\ge n_0$, and for $n\ge n_0$, $\prod_{k\ge n}r_k=a/\prod_{n_0\le k<n}r_k\to1$. By \eqref{eq:lecam},
$\dtv(P_n,Q)\le\sqrt{1-\rho(P_n,Q)^2}\to0$.
\end{proof}

\begin{example}\label{ex:oscillating}
Let $E=[0,1]$, fix $\varepsilon\in(0,1]$, and let $X_k$ be independent with densities $1+\varepsilon\sin(2\pi m_kx)$,
$m_k$ positive integers. The substitution $y=m_kx$ and periodicity give
$r_k=\int_0^1\sqrt{1+\varepsilon\sin(2\pi y)}\,dy=:r$, the same for every $k$, and $r<1$ by strict concavity of the
square root. So $\prod_{k\ge n}r_k=0$ for every $n$, and by the proof of Proposition~\ref{prop:independent} there is a
single set $S\in\cE^\N$ with $\Prob(\theta_nX\in S)=1$ for all $n$ and $\lambda^{\otimes\N}(S)=0$, where $\lambda$ is
Lebesgue measure. If $m_k\to\infty$, the shifted laws converge to $\lambda^{\otimes\N}$ on every set depending on
finitely many coordinates, but not on $S$. Indeed, for a Borel set $B\subset[0,1]^{d+1}$ the probability that
$(X_n,\dots,X_{n+d})\in B$ is the integral over $B$ of $\prod_{j\le d}(1+\varepsilon\sin(2\pi m_{n+j}x_j))$.
Expanding the product, every term other than $1$ is a constant times
$\prod_{j\in J}\sin(2\pi m_{n+j}x_j)$ for a nonempty $J$, a combination of exponentials whose frequency vectors
have norm at least $\min_{j\le d}2\pi m_{n+j}\to\infty$; by the Riemann--Lebesgue lemma for $\one_B\in
L^1([0,1]^{d+1})$ its integral over $B$ tends to $0$. An explicit choice is $S=\bigcup_\ell S_\ell$, where
$S_\ell$ is the set of sequences $x$ with $\frac1N\sum_{j<N}\sin(2\pi m_{\ell+j}x_j)\to\varepsilon/2$ as $N\to\infty$. Under
the law of $\theta_nX$ the variables $\sin(2\pi m_{n+j}x_j)$, $j\in\N$, are independent, bounded, with mean
$\int_0^1\sin(2\pi m y)(1+\varepsilon\sin(2\pi my))\,dy=\varepsilon/2$, so the strong law of large numbers gives
$\Prob(\theta_nX\in S_n)=1$. Under $\lambda^{\otimes\N}$ the variables $\sin(2\pi m_{\ell+j}x_j)$ are independent and
bounded with mean $0$, so $\lambda^{\otimes\N}(S_\ell)=0$ for every $\ell$.
\end{example}

\section{Couplings}\label{sec:coupling}

In this section $E$ is a Polish space with its Borel $\sigma$-field, $Q$ is a stationary probability measure on
$E^\N$, and $X'$ has law $Q$. A \emph{coupling} of two laws, or of two processes, is a pair of random elements on one
probability space with the given marginal laws; no other condition is imposed.

\begin{lemma}\label{lem:family}
Convergence \eqref{eq:setwise} holds if and only if for every $n$ there is a coupling $(Z_n,Z'_n)$ of
$\Prob(\theta_nX\in\cdot\,)$ and $Q$ such that $\Prob(Z_n\in A,\ Z'_n\notin A)\to0$ for every $A\in\cE^\N$.
\end{lemma}

\begin{proof}
If such couplings exist, then
\[
\Prob(\theta_nX\in A)-Q(A)=\Prob(Z_n\in A,\ Z'_n\notin A)-\Prob(Z_n\in A^c,\ Z'_n\notin A^c),
\]
and both terms tend to $0$, the second by the hypothesis applied to the complement $A^c$.
Conversely, assume \eqref{eq:setwise}. Since $E^\N$ is Polish, its Borel $\sigma$-field is countably generated, so
there are finite measurable partitions $\mathcal P_1,\mathcal P_2,\dots$ of $E^\N$, each refining the previous one,
such that $\bigcup_m\sigma(\mathcal P_m)$ generates $\cE^\N$. For fixed $m$, \eqref{eq:setwise} applied to the finitely
many sets of $\sigma(\mathcal P_m)$ shows that the total variation distance $\delta_m(n)$ between
$\Prob(\theta_nX\in\cdot\,)$ and $Q$ restricted to $\sigma(\mathcal P_m)$ tends to $0$ as $n\to\infty$. Choose
$N_1<N_2<\cdots$ with $\delta_m(n)<1/m$ for $n\ge N_m$, and let $m(n)=\max\{m:N_m\le n\}$ for $n\ge N_1$ (and
$m(n)=1$ otherwise). Then $m(n)\to\infty$ and $\delta_{m(n)}(n)\to0$. For each $n$, couple the cells of
$\mathcal P_{m(n)}$ containing $\theta_nX$ and $X'$ by a maximal coupling of their two finite distributions, so that
they differ with probability $\delta_{m(n)}(n)$. Given the two cells, draw $Z_n$ and $Z'_n$ independently from the
conditional laws of $\theta_nX$ and of $X'$ given the respective cells. Each marginal is then a mixture of the
conditional laws on the cells with the correct cell probabilities, so the marginals are correct.

Let $A\in\cE^\N$ and $\varepsilon>0$. Choose $k$ and $C\in\sigma(\mathcal P_k)$ with $Q(A\,\triangle\,C)<\varepsilon$.
Once $m(n)\ge k$, $C$ is a union of cells of $\mathcal P_{m(n)}$, so on the event that $Z_n$ and $Z'_n$ lie in the
same cell, either both lie in $C$ or neither does. Hence
\[
\Prob(Z_n\in A,\ Z'_n\notin A)\le\Prob(Z_n\in A\triangle C)+\Prob(Z'_n\in A\triangle C)+\delta_{m(n)}(n),
\]
and by \eqref{eq:setwise} the right side tends to $2Q(A\triangle C)<2\varepsilon$.
\end{proof}

A coupling of the processes $X$ and $X'$ themselves is \emph{eventwise} if \eqref{eq:eventwise} holds. Applying
\eqref{eq:eventwise} to $A$ and to $A^c$, as in the first part of the preceding proof, shows that an eventwise coupling
gives \eqref{eq:setwise}.

\begin{lemma}\label{lem:closeness}
Fix a coupling of $X$ and $X'$.
\begin{enumerate}[(i)]
\item If \eqref{eq:setwise} holds and $d(X_n,X'_n)\to0$ in probability for some metric $d$ compatible with the
topology of $E$, the coupling is eventwise.
\item If the coupling is eventwise, then $d_0(X_n,X'_n)\to0$ in probability for every separable metric $d_0$ on $E$
whose Borel sets are those of $E$.
\end{enumerate}
\end{lemma}

\begin{proof}
(i) Give $E^\N$ the metric $\bar d(x,y)=\sum_k2^{-k}\min\{1,d(x_k,y_k)\}$, which induces the product topology. For each
$k$, $d(X_{n+k},X'_{n+k})\to0$ in probability as $n\to\infty$, so $\bar d(\theta_nX,\theta_nX')\to0$ in probability.
Let $A\in\cE^\N$ and $\varepsilon>0$. Since $Q$ is a Borel probability measure on a Polish space, there are a closed
set $F$ and an open set $U$ with $F\subset A\subset U$ and $Q(U\setminus F)<\varepsilon$, and a compact $K\subset F$
with $Q(F\setminus K)<\varepsilon$. If $K$ is empty or $U=E^\N$, put $\delta=1$. Otherwise the compact set $K$ and
the nonempty closed set $E^\N\setminus U$ are disjoint, so their distance $\delta$ is positive. In all cases, if
$\theta_nX\in K$ and $\bar d(\theta_nX,\theta_nX')<\delta$, then $\theta_nX'\in U$. Hence
\[
\Prob(\theta_nX\in A,\ \theta_nX'\notin A)\le\Prob(\theta_nX\in A\setminus K)+\Prob\bigl(\bar d(\theta_nX,\theta_nX')
\ge\delta\bigr)+\Prob(\theta_nX'\in U\setminus A).
\]
By \eqref{eq:setwise} the first term tends to $Q(A\setminus K)<2\varepsilon$, the second tends to $0$, and the third
equals $Q(U\setminus A)<\varepsilon$ by stationarity.

(ii) Let $\delta>0$. By separability, $E$ is covered by countably many Borel sets $B_1,B_2,\dots$ of $d_0$-diameter
less than $\delta$. If $d_0(X_n,X'_n)\ge\delta$, then either $X_n\notin B_1\cup\dots\cup B_N$, or $X_n\in B_i$ and
$X'_n\notin B_i$ for some $i\le N$. Hence
\[
\Prob(d_0(X_n,X'_n)\ge\delta)\le\sum_{i\le N}\Prob(X_n\in B_i,\ X'_n\notin B_i)+\Prob(X_n\notin B_1\cup\dots\cup B_N).
\]
Each term of the sum tends to $0$ by \eqref{eq:eventwise} applied to $\{x_0\in B_i\}$, and the last term tends to
$Q(x_0\notin B_1\cup\dots\cup B_N)$ by \eqref{eq:setwise}, which is small for large $N$.
\end{proof}

The next proposition shows that eventwise couplings exist in the setting of Theorem~\ref{thm:main} when $T$ is a shift
whose invariant measure is trivial on the right tail $\bigcap_m\sigma(y_k:k\ge m)$ of the coordinates $y_k$, as for the Bernoulli shift of Example~\ref{ex:bernoulli}. The proof uses
Goldstein's maximal coupling theorem \cite{Gol79}; see also \cite[Chapter~4]{Tho00}. It states the following. Let $Z$
and $Z'$ be discrete-time processes with values in a Polish space such that
\[
\dtv\bigl(\Prob(\theta_nZ\in\cdot\,),\Prob(\theta_nZ'\in\cdot\,)\bigr)\to0 .
\]
Then there are a coupling of $Z$ and $Z'$ and an almost surely finite random time $M$ with $Z_k=Z'_k$ for all $k\ge M$.

\begin{proposition}\label{prop:Ktail}
Let $S$ be a Polish space, $\Omega=S^\Z$, $T$ the shift $(Ty)_j=y_{j+1}$, and $\nu$ a shift-invariant probability
measure on $\Omega$ that is trivial on the right tail $\mathcal R=\bigcap_m\mathcal G_m$, where
$\mathcal G_m=\sigma(y_k:k\ge m)$. Let $f$, $X$ and $X'$ be as in Theorem~\ref{thm:main}. Then \eqref{eq:setwise} holds,
$\dtv(\Prob(\theta_nX\in\cdot\,),\Prob(X'\in\cdot\,))=\frac12\int|f-1|\,d\nu$ for every $n$, and $X$ and $X'$ admit an
eventwise coupling.
\end{proposition}

\begin{proof}
By reverse martingale convergence and triviality of $\mathcal R$, $\E_\nu[f\mid\mathcal G_m]\to\E_\nu[f\mid\mathcal
R]=1$ in $L^1(\nu)$.

\emph{Convergence \eqref{eq:setwise}.} Let $B\subset\Omega$ be a Borel set depending only on the coordinates in
$[-w,w]$. Then $\one_B\circ T^n$ is $\mathcal G_{n-w}$-measurable, so
\[
\bigl|\Prob(T^nY\in B)-\nu(B)\bigr|=\Bigl|\int(\one_B\circ T^n)\bigl(\E_\nu[f\mid\mathcal G_{n-w}]-1\bigr)\,d\nu\Bigr|
\le\bigl\|\E_\nu[f\mid\mathcal G_{n-w}]-1\bigr\|_{L^1(\nu)}\to0 .
\]
For a general Borel set $B$ and $\varepsilon>0$, choose $\eta\in(0,\varepsilon]$ such that $\int_Df\,d\nu<\varepsilon$
whenever $\nu(D)<\eta$, and a set $C$ depending on finitely many coordinates with $\nu(B\triangle C)<\eta$. By invariance
$\nu(T^{-n}(B\triangle C))<\eta$, so $\Prob(T^nY\in B\triangle C)<\varepsilon$ for every $n$, and
$\limsup_n|\Prob(T^nY\in B)-\nu(B)|\le2\varepsilon$. Since every path event for $X$ is of the form
$\{\theta_nX\in A\}=\{T^nY\in\Psi^{-1}A\}$ with $\Psi$ as in the proof of Theorem~\ref{thm:main}, \eqref{eq:setwise}
follows. The total variation identity holds for every invertible measure-preserving $T$, as in the proof of
Theorem~\ref{thm:main}.

\emph{The coupling.} The law of $(Y_{n+k})_{k\ge0}$ is the restriction of $f\,d\nu$ to $\mathcal G_n$, transported to
$S^\N$, and its density with respect to the corresponding restriction of $\nu$ is $\E_\nu[f\mid\mathcal G_n]$. Hence
the total variation distance between the laws of $(Y_{n+k})_{k\ge0}$ and $(Y'_{n+k})_{k\ge0}$ equals
$\frac12\|\E_\nu[f\mid\mathcal G_n]-1\|_{L^1(\nu)}\to0$. By Goldstein's theorem there is a coupling of the one-sided
sequences $(Y_k)_{k\ge0}$ and $(Y'_k)_{k\ge0}$ with $Y_k=Y'_k$ for all $k\ge M$, $M$ almost surely finite. Extend it to
the two-sided sequences by drawing $(Y_k)_{k<0}$ and $(Y'_k)_{k<0}$ independently from the regular conditional laws
given the respective right halves, under $f\,d\nu$ and under $\nu$; these exist because $S$ is Polish. This gives a
coupling of $Y$ and $Y'$, hence of $X=\Psi(Y)$ and $X'=\Psi(Y')$.

On $\Omega$ use the metric $d(y,z)=\sum_{k\in\Z}2^{-|k|}\min\{1,d_S(y_k,z_k)\}$, where $d_S$ is a complete metric for
$S$; $\Omega$ is Polish. The points $X_n=T^nY$ and $X'_n=T^nY'$ have the same coordinates of index $k\ge M-n$, so
$d(X_n,X'_n)\le\sum_{k<M-n}2^{-|k|}\to0$ almost surely. Lemma~\ref{lem:closeness}(i) shows that the coupling is
eventwise.
\end{proof}

For the Bernoulli shift of Example~\ref{ex:bernoulli} the coupling is explicit: take $Y'$ fair and let $Y$ equal $Y'$ with the coordinate $0$
replaced by $1$.

\begin{proof}[Proof of Theorem~\ref{thm:horocycle}]
Write $\phi_t(x)=xu_t$ for the horocycle flow on $\Omega$, so that $T=\phi_1$. The flow is mixing:
$\nu(\phi_t^{-1}C\cap D)\to\nu(C)\nu(D)$ as $t\to\infty$ \cite[Remark~3.2.8]{Mor05}. In particular the invertible map
$T$ is mixing, and Theorem~\ref{thm:main} gives \eqref{eq:setwise}.

Suppose that some coupling of $X$ and $X'$ has $d(X_n,X'_n)\to0$ in probability. Under the law of $X$, the set
$\{x:x_n=T^nx_0\}$ (measurable, since $\Omega$ is a separable metric space) has probability one, and likewise for $X'$.
Put $Y=X_0$, $Y'=X'_0$ and $Z=(Y,Y')$, a random point of $\Omega\times\Omega$ with law $\lambda$, say; then $Y$ has law
$f\,d\nu$, $Y'$ has law $\nu$, and $d(T^nY,T^nY')\to0$ in probability. The map $(s,x)\mapsto\phi_s(x)$ is uniformly
continuous on the compact set $[0,1]\times\Omega$, so for every $\varepsilon>0$ there is $\eta>0$ such that
$d(x,y)<\eta$ implies $d(\phi_sx,\phi_sy)<\varepsilon$ for all $s\in[0,1]$. Hence $d(\phi_tY,\phi_tY')\to0$ in
probability as $t\to\infty$ through the reals. Let $q(y,y')=\min\{1,d(y,y')\}$ and
$\Phi_t=\phi_t\times\phi_t$. Put $W_t=q(\Phi_tZ)$. Then $0\le W_t\le1$ and $W_t\to0$ in probability, so for every $\varepsilon>0$
\[
\E\,W_t\le\varepsilon+\Prob(W_t>\varepsilon)\to\varepsilon ,
\]
hence $\E\,W_t\to0$ and therefore
\begin{equation}\label{eq:avg}
\int\Bigl(\frac1S\int_0^Sq(\Phi_tz)\,dt\Bigr)\lambda(dz)=\frac1S\int_0^S\E\,q(\Phi_tZ)\,dt\to0\qquad(S\to\infty).
\end{equation}

The space $\Omega\times\Omega$ is the quotient of $G\times G$ by the lattice $\Gamma\times\Gamma$, and $\Phi_t$ is right
multiplication by the one-parameter unipotent subgroup $\{(u_t,u_t)\}$ of $G\times G$. By Ratner's equidistribution
theorem \cite{Rat91}, in the form of \cite[Theorem~1.3.4 and Remark~1.3.5]{Mor05}, for every $z\in\Omega\times\Omega$
there is a closed connected subgroup $H_z$ of $G\times G$ containing $\{(u_t,u_t)\}$ such that the orbit $zH_z$ is
closed and carries an $H_z$-invariant probability measure $\mu_z$. The $\Phi$-orbit of $z$ is dense in $zH_z$, and
$\frac1S\int_0^Sg(\Phi_tz)\,dt\to\int g\,d\mu_z$ for every continuous $g$. The support of $\mu_z$ is $zH_z$, which
contains $z$. The function
$a(z)=\int q\,d\mu_z=\lim_{S\to\infty}\frac1S\int_0^Sq(\Phi_tz)\,dt$ is a pointwise limit of continuous functions,
hence measurable, and by \eqref{eq:avg} and bounded convergence $\int a\,d\lambda=0$. So for $\lambda$-almost every $z$
the closed set $\{q=0\}$, the diagonal of $\Omega\times\Omega$, has full $\mu_z$-measure and therefore contains the
support of $\mu_z$, and so contains $z$. Thus $Y=Y'$ almost surely, which forces $f\,d\nu=\nu$, contrary to the
assumption on $f$.

Finally, an eventwise coupling would give $d(X_n,X'_n)\to0$ in probability by Lemma~\ref{lem:closeness}(ii).
\end{proof}

\begin{remark}\label{rem:problem33}
Let $\bar d(x,y)=\sum_k2^{-k}\min\{1,d(x_k,y_k)\}$ be the product metric on $\Omega^\N$ of
Lemma~\ref{lem:closeness}. Since $\min\{1,d(X_n,X'_n)\}\le\bar d(\theta_nX,\theta_nX')$, Theorem~\ref{thm:horocycle} also shows that no coupling of $X$ and $X'$ satisfies
$\bar d(\theta_nX,\theta_nX')\to0$ in probability, although $\theta_nX$ converges to $X'$ setwise, and therefore in
distribution.
\end{remark}

\begin{remark}\label{rem:open}
Eventwise couplings are sufficient for
\eqref{eq:setwise} but, by Theorem~\ref{thm:horocycle}, not necessary. A characterization would have to allow some
other form of single coupling, for instance with a random time change, as shift-coupling does for Ces\`aro
convergence. It would also be of interest to determine the stationary mixing laws $Q$ for which an eventwise coupling
exists for every density $h$; Proposition~\ref{prop:Ktail} gives shifts with a trivial right tail, and
Theorem~\ref{thm:horocycle} excludes the horocycle flow.
\end{remark}

\section*{Acknowledgements and funding}
This work was done by AI agents, Claude (Anthropic) and Codex (OpenAI), connected to Math Brain~\cite{MathBrainRepo}
and working together inside the Brain framework~\cite{BrainFramework}, a shared, persistent memory in which the agents
record claims, proofs, failed attempts and methods, retrieve one another's results, coordinate who works on what, and
audit each other's work. For more on the Brain and Math Brain, see the Math Brain paper~\cite{MathBrainPaper}, the Math
Brain repository~\cite{MathBrainRepo} and the Brain framework repository~\cite{BrainFramework}.

The research presented in this paper was supported by the European Research Council (ERC) under the European Union's
Horizon Europe research and innovation programme (grant agreement No.~101041711), by the Simons Foundation as part of
the Collaboration on the Mathematical and Scientific Foundations of Deep Learning, by Heights Labs, by Convex Nexus
Capital, by the Israel Science Foundation (grant number 2258/19), and by the Israel Science Foundation (ISF Grant
4101/25).

\section*{Declaration of competing interest}
The author holds equity in Heights Labs and in Convex Nexus Capital and has a management role at Convex Nexus Capital.
These entities are included in the funding declaration above.

\section*{Data availability}
No data were generated or analysed in this work.

\section*{Declaration of generative AI and AI-assisted technologies}
This paper is part of a research project on AI proving problems. Claude agents (Anthropic) found the constructions and
proofs and wrote the manuscript. Codex agents (OpenAI) checked the
proofs and the references independently. The results used are due to Thorisson \cite{Tho00,Tho11}, Lin \cite{Lin71}, Kakutani \cite{Kak48}, Goldstein \cite{Gol79}, Ratner \cite{Rat91}, and Creutz, Pavlov and Rodock \cite{CPR23}. The author
directed the work and is responsible for the content of the paper
and for decisions concerning its dissemination.

\begin{thebibliography}{99}
\bibitem{CPR23} D.~Creutz, R.~Pavlov and S.~Rodock, \emph{Measure-theoretically mixing subshifts with low complexity},
Ergodic Theory Dynam. Systems \textbf{43} (2023), no.~7, 2293--2316, \doi{10.1017/etds.2022.42}; \arxiv{2201.00489}.
\bibitem{Gol79} S.~Goldstein, \emph{Maximal coupling}, Z. Wahrsch. Verw. Gebiete \textbf{46} (1979), no.~2, 193--204,
\doi{10.1007/BF00533259}.
\bibitem{Kak48} S.~Kakutani, \emph{On equivalence of infinite product measures}, Ann. of Math. (2) \textbf{49} (1948),
no.~1, 214--224, \doi{10.2307/1969123}.
\bibitem{Lin71} M.~Lin, \emph{Mixing for Markov operators}, Z. Wahrsch. Verw. Gebiete \textbf{19} (1971), no.~3,
231--242, \doi{10.1007/BF00534111}.
\bibitem{Mor05} D.~W. Morris, \emph{Ratner's Theorems on Unipotent Flows}, Chicago Lectures in Mathematics, University
of Chicago Press, Chicago, 2005; \arxiv{math/0310402}.
\bibitem{Rat91} M.~Ratner, \emph{Raghunathan's topological conjecture and distributions of unipotent flows}, Duke Math.
J. \textbf{63} (1991), no.~1, 235--280, \doi{10.1215/S0012-7094-91-06311-8}.
\bibitem{Tho00} H.~Thorisson, \emph{Coupling, Stationarity, and Regeneration}, Probability and its Applications,
Springer, New York, 2000, \doi{10.1007/978-1-4612-1236-2}.
\bibitem{Tho11} H.~Thorisson, \emph{Open problems in renewal, coupling and Palm theory}, Queueing Syst. \textbf{68}
(2011), no.~3--4, 313--319, \doi{10.1007/s11134-011-9241-2}.
\bibitem{ThoPre} H.~Thorisson, \emph{Some open probability problems}, preprint version of \cite{Tho11},
\url{https://web.archive.org/web/20240327021031id_/http://cms.dm.uba.ar/depto/public/Some\%20Open\%20Problems-preprint.pdf}.

\bibitem{WS26} We-Solve, \emph{A sharp obstruction to two-process asymptotic coupling under shifted weak convergence},
model-assisted contribution Math-11651-A1, we-solve, 28 October 2026,
\url{https://we-solve.github.io/contributions/two-process-coupling-characterization-of-weak-convergence/attempt-1/r1/}.
\bibitem{Zer26} R.~Zeraoulia, \emph{Setwise asymptotic stationarity does not imply total variation asymptotic
stationarity: an explicit counterexample}, SSRN preprint 6361341, 2026, \doi{10.2139/ssrn.6361341},
\url{https://papers.ssrn.com/sol3/papers.cfm?abstract_id=6361341}.
\bibitem{MathBrainRepo} Y.~Shmalo, \emph{Math Brain: public snapshot, code and tests}, 2026,
\url{https://github.com/yspennstate/math-brain-public/tree/8c9a6bd478c4ec4082347151b7ae16db50c917dd}.
\bibitem{MathBrainPaper} Y.~Shmalo, \emph{Math Brain: shared, auditable memory for multi-agent mathematical
research}, 2026, \doi{10.5281/zenodo.24150837},
\url{https://github.com/yspennstate/math-brain-public/tree/main/papers/math-brain-2026-10-24}.
\bibitem{BrainFramework} Y.~Shmalo, \emph{Brain Framework}, public software repository, 2026,
\url{https://github.com/yspennstate/brain-framework-public/tree/5629e644cdd510173ae97e9aa92642dabd22dc57}.
\end{thebibliography}

\end{document}
